\begin{theorem}[Ловас, 1972; слабая теорема о совершенных графах]
Граф \(G\) совершенен тогда и только тогда, когда \(\overline G\) совершенен.
\end{theorem}
Достаточно доказать «\(\Rightarrow\)» (обратное — то же для \(\overline G\), так как \(\overline{\overline G}=G\)).
Нам понадобятся три леммы.

\begin{lemma}[\(\alpha\chi\ge n\)]\label{lem:alphachi}
Для любого графа \(\alpha(G)\,\chi(G)\ge n\).
\end{lemma}
\begin{proof}
Покрасим \(G\) в \(\chi(G)\) цветов; классы цветов \(V_1,\ldots,V_{\chi}\) независимы, поэтому
\(n=\sum|V_i|\le\chi(G)\,\alpha(G)\).
\end{proof}

\begin{definition}[Клонирование вершины]
Для \(x\in V\) граф \(G\circ x\) получается добавлением новой вершины \(x'\), смежной с \(x\) и
со всеми соседями \(x\): \(G\circ x=\bigl(V\cup\{x'\},\;E\cup\{x'u: xu\in E\}\cup\{xx'\}\bigr)\).
\end{definition}
\begin{lemma}[Клонирование сохраняет совершенность]\label{lem:clone}
Если \(G\) совершенен, то \(G\circ x\) совершенен.
\end{lemma}
\begin{proof}
Индуцированный подграф \(G\circ x\) — либо индуцированный подграф \(G\) (если в нём нет
\(x\) или \(x'\); вершина \(x'\) без \(x\) играет роль \(x\)), либо \(H\circ x\) для индуцированного
\(H\le G\), \(x\in H\), который тоже совершенен. Поэтому достаточно доказать
\(\chi(G\circ x)\le\omega(G\circ x)\) для совершенного \(G\); пусть \(\omega=\omega(G)\).

\emph{1) \(x\) лежит в некоторой максимальной клике \(K\).} Тогда \(K\cup\{x'\}\) — клика, и
\(\omega(G\circ x)=\omega+1\). Покрасим \(G\) в \(\chi(G)=\omega\) цветов, а \(x'\) — в новый цвет:
\(\chi(G\circ x)\le\omega+1\).

\emph{2) \(x\) не лежит ни в одной максимальной клике.} Клика с \(x\) и \(x'\) имеет размер
\(|K|+1\le\omega\) для клики \(K\ni x\) графа \(G\), поэтому \(\omega(G\circ x)=\omega\). Покрасим \(G\) в
\(\omega\) цветов, пусть \(A\) — класс цвета \(x\). Каждая максимальная клика \(G\) содержит
вершины всех \(\omega\) цветов, значит, пересекает \(A\), причём не в \(x\) — т.\,е. пересекает
\(A\setminus\{x\}\). Поэтому \(\omega\bigl(G-(A\setminus\{x\})\bigr)\le\omega-1\), и по совершенности этот граф
красится в \(\omega-1\) цвет. Множество \((A\setminus\{x\})\cup\{x'\}\) независимо (\(x'\) смежна
только с \(x\) и \(N(x)\), а \(A\) — класс цвета \(x\)); дадим ему ещё один цвет. Итого \(\omega\) цветов.
\end{proof}
\begin{graphillustration}{Клонирование: $x'$ смежна с $x$ и со всеми соседями $x$ (штриховые рёбра). Справа случай 2: $A$ — класс цвета $x$; множество $(A\setminus\{x\})\cup\{x'\}$ независимо и получает отдельный цвет.}{fig:14-clone}
\begin{tikzpicture}[font=\small]
\begin{scope}
\node[gi vertex] (x) at (0,0) {$x$};
\node[gi vertex] (p) at (-1.2,1.2) {};
\node[gi vertex] (q) at (1.2,1.2) {};
\node[gi vertex] (r) at (0,-1.3) {};
\node[gi selected] (xx) at (1.4,-0.2) {$x'$};
\draw[gi edge] (x) -- (p); \draw[gi edge] (x) -- (q); \draw[gi edge] (x) -- (r); \draw[gi edge] (p) -- (q);
\draw[gi alternate] (xx) -- (x); \draw[gi alternate] (xx) -- (p); \draw[gi alternate] (xx) -- (q); \draw[gi alternate] (xx) -- (r);
\node[gi note] at (0,-2.0) {$G\circ x$};
\end{scope}
\begin{scope}[shift={(5.5,0)}]
\node[gi vertex] (x) at (0,0) {$x$};
\node[gi vertex] (a1) at (-1.6,-1.0) {};
\node[gi vertex] (a2) at (-0.5,-1.4) {};
\node[gi selected] (xx) at (1.4,-0.2) {$x'$};
\node[gi vertex] (n1) at (0,1.2) {};
\draw[gi edge] (x) -- (n1); \draw[gi alternate] (xx) -- (x); \draw[gi alternate] (xx) -- (n1);
\draw[gi edge] (a1) -- (n1); \draw[gi edge] (a2) -- (n1);
\node[gi group,fit=(x)(a1)(a2),label=left:{$A$}] {};
\node[gi note,anchor=west] at (2.2,0.2) {новый цвет:\\$(A\setminus\{x\})\cup\{x'\}$};
\end{scope}
\end{tikzpicture}
\end{graphillustration}


\begin{lemma}[Критерий через независимое множество]\label{lem:crit}
\(G\) совершенен \(\iff\) в каждом индуцированном \(H\le G\) есть независимое множество,
пересекающее все максимальные (размера \(\omega(H)\)) клики \(H\).
\end{lemma}
\begin{proof}
\(\Rightarrow\). Покрасим \(H\) в \(\chi(H)=\omega(H)\) цветов. Максимальная клика содержит вершины
всех цветов, поэтому класс любого цвета — искомое независимое множество.

\(\Leftarrow\). Докажем \(\chi(H)\le\omega(H)\) для всех \(H\le G\) индукцией по \(\omega(H)\) (лекционный
«алгоритм»: красим \(S\) в цвет \(i\), удаляем, \(i\gets i+1\)). При \(\omega(H)=0\) граф пуст. Пусть
\(S\) — независимое множество из условия. Каждая клика размера \(\omega(H)\) теряет вершину, поэтому
\(\omega(H-S)=\omega(H)-1\). Граф \(H-S\le G\) по индукции красится в \(\omega(H)-1\) цветов; \(S\) получает
ещё один.
\end{proof}

\begin{proof}[Доказательство теоремы Ловаса]
Пусть \(G\) совершенен. Клики \(\overline G\) — независимые множества \(G\), и наоборот;
индуцированные подграфы \(\overline G\) — дополнения индуцированных подграфов \(G\), а они
совершенны. По лемме~\ref{lem:crit}, применённой к \(\overline G\), достаточно доказать:
\begin{center}
\emph{в любом совершенном графе \(G\) есть клика \(K\), пересекающая все максимальные
независимые множества.}
\end{center}
Пусть \(\mathbb A\) — множество всех максимальных (размера \(\alpha=\alpha(G)\)) независимых
множеств, \(k=|\mathbb A|\), и для вершины \(v\)
\[
a_v=\bigl|\{I\in\mathbb A: v\in I\}\bigr| .
\]
Построим \(G'\): каждую вершину \(v\) заменим \(a_v\) попарно смежными копиями (клонируем
\(a_v-1\) раз; при \(a_v=0\) удаляем \(v\)). По лемме~\ref{lem:clone} и наследственности \(G'\) совершенен.

\emph{Число вершин.} \(v(G')=\sum_v a_v=\sum_v\sum_{I\in\mathbb A}[v\in I]=\sum_{I\in\mathbb A}|I|=k\alpha\).

\emph{\(\alpha(G')=\alpha\).} Независимое множество \(G'\) содержит не более одной копии каждой
вершины, и исходные вершины образуют независимое множество \(G\).

\emph{\(\chi(G')\le k\).} Распределим \(a_v\) копий \(v\) по одной между \(a_v\) множествами \(I\ni v\).
Копии, отнесённые к \(I\), образуют независимое множество (копия \(I\)); это раскраска
в \(k\) цветов, где цвет \(I\) состоит из копий вершин \(I\).

Из леммы~\ref{lem:alphachi}: \(k\alpha=v(G')\le\alpha(G')\chi(G')\le\alpha k\), значит,
\(\chi(G')=k\), и по совершенности \(\omega(G')=k\). Пусть \(K'\) — клика \(G'\) размера \(k\).
В раскраске выше \(k\) цветов, а вершины \(K'\) окрашены по-разному, поэтому \(K'\) содержит
ровно одну вершину каждого цвета.

Пусть \(K\) — множество вершин \(G\), копии которых лежат в \(K'\). Это клика (копии разных
вершин смежны, только если смежны оригиналы). Для любого \(I\in\mathbb A\) клика \(K'\)
содержит вершину цвета \(I\) — копию некоторой \(v\in I\); тогда \(v\in K\cap I\).
Итак, \(K\) пересекает все максимальные независимые множества.
\end{proof}
\begin{graphillustration}{Конструкция Ловаса для $G=P_4=abcd$. Максимальные независимые множества: $\{a,c\},\{a,d\},\{b,d\}$, $k=3$, $\alpha=2$; $a_a=a_d=2$, $a_b=a_c=1$. В $G'$ ровно $k\alpha=6$ вершин, клика $\{a_1,a_2,b\}$ размера $k$ проецируется в клику $K=\{a,b\}$, пересекающую все три множества.}{fig:14-blowup}
\begin{tikzpicture}[font=\small]
\begin{scope}
\node[gi vertex] (a) at (0,0) {$a$};
\node[gi vertex] (b) at (1.2,0) {$b$};
\node[gi vertex] (c) at (2.4,0) {$c$};
\node[gi vertex] (d) at (3.6,0) {$d$};
\draw[gi edge] (a) -- (b) -- (c) -- (d);
\node[gi note] at (1.8,-0.8) {$G=P_4$};
\end{scope}
\node[gi note] at (4.8,0) {$\Longrightarrow$};
\begin{scope}[shift={(6.2,0)}]
\node[gi selected] (a1) at (0,0.7) {$a_1$};
\node[gi selected] (a2) at (0,-0.7) {$a_2$};
\node[gi selected] (b) at (1.4,0) {$b$};
\node[gi vertex] (c) at (2.8,0) {$c$};
\node[gi vertex] (d1) at (4.2,0.7) {$d_1$};
\node[gi vertex] (d2) at (4.2,-0.7) {$d_2$};
\draw[gi match] (a1) -- (a2); \draw[gi match] (a1) -- (b); \draw[gi match] (a2) -- (b);
\draw[gi edge] (b) -- (c); \draw[gi edge] (c) -- (d1); \draw[gi edge] (c) -- (d2); \draw[gi edge] (d1) -- (d2);
\node[gi note] at (2.1,-1.4) {$G'$: $v(G')=6$, $\chi(G')=\omega(G')=3$};
\end{scope}
\end{tikzpicture}
\end{graphillustration}


\begin{tcolorbox}[colback=cbGrayBack,colframe=cbGray!55]
\textbf{Что ещё полезно знать (без доказательств).}
\begin{itemize}\tightlist
\item \emph{Критерий Ловаса:} \(G\) совершенен \(\iff\) \(\alpha(H)\,\omega(H)\ge v(H)\) для всех \(H\le G\);
  он симметричен относительно дополнения.
\item \emph{Сильная теорема (Чудновская–Робертсон–Сеймур–Томас, 2002; гипотеза Бержа):}
  \(G\) совершенен \(\iff\) ни \(G\), ни \(\overline G\) не содержат индуцированного нечётного цикла
  длины \(\ge5\) («нечётных дыр и антидыр»).
\item \emph{Хордальные графы} (нет индуцированных циклов длины \(\ge4\)) совершенны; это
  обобщает интервальные графы (билет~\ref{ticket:13}).
\item \emph{Следствие (Дилуорс):} дополнение графа сравнимости совершенно.
\end{itemize}
\end{tcolorbox}
